\definecolor{BleuLogo}{RGB}{79,129,150}
\definecolor{MarronLogo}{RGB}{210,163,128}
\definecolor{RoseLogo}{RGB}{250,207,187}

\thispagestyle{empty}

\begin{tikzpicture}[remember picture,overlay,
        every shadow/.style={opacity=.6,fill=black!50,shadow xshift=5pt, shadow yshift=-5pt}]

% vague colorée
    \fill[color=BleuLogo!25]
        ($(current page.south west)$)..controls+(0:20)and+(-160:5)..($(current page.east)$)--
        ($(current page.south east)$)--cycle;

% courbe colorée sur la vague
    \draw[line width=3pt,color=BleuLogo!100,shorten <=-5pt]
        ($(current page.south west)$)..controls+(0:20)and+(-160:5)..($(current page.east)$)
        node[pos=0.75] (c) {\tikz{\fill(0,0) circle (.2cm);}}; % pour le point

% titre
    \node[anchor=north east,font=\itshape\Huge,color=MarronLogo,inner sep=10pt,scale=3]
        at ($(current page.north east)$) {{\hminfamily Initiation}};

% Logo JPT
    \node[anchor=north west,inner sep=4pt,scale=0.75] (logo)
        at ($(current page.north west)+(1cm,-0.5cm)$)
            {\begin{tabular}{c}
                \includegraphics[width=4.5cm]{JPT} \\
                \makebox[4.5cm][s]{\bfseries Lycée Jean-Pierre \bsc{Timbaud}}
            \end{tabular}};

% Auteur
    \node[anchor=south east,inner sep=10pt,scale=1.5,color=black!100] (auteur)
        at ($(current page.south east)$)
        {\bfseries Philippe \bsc{De Sousa}};

% 2e Courbe colorée + node avec cadre exemple maths
    \draw[line width=3pt,color=BleuLogo,shorten <=-5pt]
        (c)..controls+(110:30)and+(-70:23)..($(current page.north west)+(0,-0.1\paperheight)$)
        node[pos=0.9,color=black,fill=white,draw=RoseLogo!50,rounded corners=15pt,
                    inner sep=10pt,drop shadow,anchor=west,scale=1.5]
            {
              $\int_a^b f(x)dx = F(b) - F(a)$
            };

% LaTeX
    \node[scale=30,color=black!15,rotate=45,opacity=0.65] (latex)
        at ($(current page.center)+(1cm,0)$) {\LaTeX};

% cadre Partie 1
    \node[line width=3pt,color=black,fill=white,draw=RoseLogo!50,rounded corners=15pt,
                inner sep=10pt,drop shadow,anchor=east,scale=2.5]
        at (latex.south)
            {%
              \bfseries\scshape 1\up{\upshape ère} Partie
            } ;

\end{tikzpicture} 